% Generated draft: requires mathematical and source review.
\documentclass[11pt]{article}
\usepackage[margin=25mm]{geometry}
\usepackage{fontspec}
\setmainfont{Latin Modern Roman}
\usepackage{amsmath,amssymb,mathtools,unicode-math}
\setmathfont{Latin Modern Math}
\usepackage{enumitem,needspace,xurl,hyperref}
\setlength{\parindent}{0pt}
\setlength{\parskip}{5pt}
\setlength{\emergencystretch}{3em}
\newcommand{\N}{\mathbb{N}}
\newcommand{\Z}{\mathbb{Z}}
\newcommand{\Q}{\mathbb{Q}}
\newcommand{\R}{\mathbb{R}}
\newcommand{\C}{\mathbb{C}}
\newcommand{\F}{\mathbb{F}}
\newcommand{\A}{\mathbb{A}}
\newcommand{\PP}{\mathbb{P}}
\newcommand{\E}{\mathbb{E}}
\newcommand{\eps}{\varepsilon}
\newcommand{\dd}{\,\mathrm{d}}
\begin{document}
\section*{UnsolvedMath ID: 30002614}
\textbf{Source problem code:} catalogue entry
\subsection{Definitions and mathematical statement}
Let \(M_{15,16}\) denote the space described in the supplied record as the Brill--Noether moduli space of curves of genus \(15\) and degree \(16\). This description does not specify the moduli functor: degree requires additional data, such as a line bundle or a projective embedding. The record does not specify the dimension of a linear series, its completeness or base-point conditions, or the equivalence relation on the parametrized objects. Thus \(M_{15,16}\) is retained as source notation, not identified here with a particular moduli space.

For the usual characteristic-zero interpretation, let \(X\) be a smooth projective integral variety over an algebraically closed field of characteristic zero. A variety is rationally connected if two general geometric points lie on a rational curve. A maximal rationally connected fibration of \(X\) is a dominant rational map \(\pi\colon X\dashrightarrow B\), with connected general fibers, whose general fiber is rationally connected and whose general fibers are the equivalence classes generated by chains of rational curves through general points of \(X\). Its base is considered up to birational equivalence, so \(\dim B\) is the relevant invariant.

\textbf{Source question, conditional on specifying the moduli space and the applicable framework.} If \(X\) is a smooth projective model of the intended irreducible space \(M_{15,16}\), and \(\pi\colon X\dashrightarrow B\) is its maximal rationally connected fibration, is
\[
\dim B=3?
\]
The characteristic-zero framework above explains the terminology; it is not a field hypothesis supplied by the source record. The ground field, precise moduli problem, component, and treatment of compactification remain unresolved. Consequently, this is a conditional transcription of the intended question, not a fully specified conjecture about a uniquely defined moduli space.
\subsection{Short English statement}
The question asks whether the maximal rationally connected quotient of the intended genus-\(15\), degree-\(16\) Brill--Noether moduli space has dimension \(3\). The record leaves the moduli data and ground field unspecified. Its open status and literature assessment are supplied dated curation, requiring independent review.
\subsection{Sources}
\begin{itemize}
\item Local UnsolvedMath registry, supplied record with permanent ID \(30002614\), exact source problem code \texttt{catalogue entry}, title: ``Rationally Connected Fibration of a Brill--Noether Moduli Space.'' The supplied set is Oberwolfach Reports Open Problems; the proposed year is \(2014\). No individual proposer is identified.
\item Supplied report citation: \emph{Classical Algebraic Geometry (2015)}, Oberwolfach Reports, DOI 10.4171/owr/2014/31, pp.\ 1695--1745. \url{https://doi.org/10.4171/owr/2014/31}. The background describes these reports as records of mathematical workshops and their open-problem discussions.
\item Supplied literature reference: ``Schreyer, matrix factorizations and genus 15.'' \url{https://doi.org/10.14231/AG-2015-021}.
\item Supplied literature reference: ``Schreyer, original preprint.'' \url{https://arxiv.org/abs/1311.6962}.
\item Supplied literature reference: ``Oberwolfach Report 31/2014.'' \url{https://ems.press/content/serial-article-files/46521}.
\item Literature assessment dated exactly \texttt{2026-08-21}: the supplied curation labels the problem open and reports that Schreyer's published paper presents the three-dimensional-base statement only as evidence for a conjecture about the maximal rationally connected fibration of the Brill--Noether moduli space; the curation reports finding no proof of that conjecture. This is supplied dated curation derived from the OWR record, not an independently verified literature assessment. The background explicitly requests independent verification of time-sensitive status and cited claims.
\item The registry labels the statement \texttt{corrected\_verified} and records a manual semantic, typographic, and LaTeX review asserting a meaning-preserving cleanup. These are supplied editorial claims, not independent verification of the mathematical formulation or cited sources.
\end{itemize}
\end{document}
